\begin{question}[type=true_false,difficulty=4,code={II-4}]
Một doanh nghiệp kinh doanh sản xuất đồng hồ có đồ thị hàm tổng chi phí theo số sản phẩm, là một phần của đồ thị hàm số $ f(x)=\frac{ax^2+bx+c}{x+e} $ như hình vẽ (mỗi đơn vị trên trục hoành tương ứng với 100 sản phẩm, mỗi đơn vị trên trục tung tương ứng với 1000 USD). Biết rằng tầm đối xứng của đồ thị hàm số đó là điểm $ I\left(-1;\frac{2}{3}\right) $ và đường tiệm cận xiên của đồ thị đó đi qua điểm $ B(3;2) $.

\image[alt={Hình minh hoạ}]{https://pending-r2.local/raw/16_ee6adab79a.jpg}


\begin{statement}[correct=false,label=a]
Hàm số đồng biến trên khoảng $(0;+\infty)$.
\end{statement}

\begin{statement}[correct=false,label=b]
Đồ thị hàm số có đường tiệm cận đứng là $x = 1$.
\end{statement}

\begin{statement}[correct=true,label=c]
Hàm số có thể viết lại dưới dạng $f(x)=\frac{1}{3}x+1+\frac{d}{x+1}$, với $d \in \mathbb{R}$.
\end{statement}

\begin{statement}[correct=false,label=d]
Theo khảo sát, tổng doanh thu được mô tả bằng hàm số $R(x) = x^{2} + 2x$ và lợi nhuận khi bán 200 sản phẩm là 5250 USD. Khi chi phí nhỏ nhất, số sản phẩm (làm tròn) là 25 sản phẩm.
\end{statement}

\end{question}